\section{Complete vector and context catalogue}
\label{sec:complete-catalogue}
Each ray is projective: any nonzero scalar multiple denotes the same ray. The ordering below is exactly the ordering stored in \path{audit_v3/vectors.json}.  Context indices are one-based in this catalogue.
\subsection{Firefly logic \texorpdfstring{$L_{12}$}{L12}}
\begingroup\sloppy Source: D. W. Cohen, An Introduction to Hilbert Space and Quantum Logic, Springer (1989), DOI 10.1007/978-1-4613-8841-8.  Construction: explicit five-ray firefly realization. The audit verifies 5 projectively distinct rays and 2 complete contexts.\par\endgroup
\begin{longtable}{@{}r l@{}}
\toprule
ray & representative\\
\midrule
\endfirsthead
\toprule
ray & representative\\
\midrule
\endhead
$v_{1}$ & $(1,0,0)$\\
$v_{2}$ & $(0,1,0)$\\
$v_{3}$ & $(0,0,1)$\\
$v_{4}$ & $(0,1,1)$\\
$v_{5}$ & $(0,1,-1)$\\
\bottomrule
\end{longtable}
\begin{longtable}{@{}r l@{}}
\toprule
context & mutually orthogonal rays\\
\midrule
\endfirsthead
\toprule
context & mutually orthogonal rays\\
\midrule
\endhead
$C_{1}$ & $\{v_{1},v_{2},v_{3}\}$\\
$C_{2}$ & $\{v_{1},v_{4},v_{5}\}$\\
\bottomrule
\end{longtable}
\subsection{Yu--Oh 13}
\begingroup\sloppy Source: S. Yu and C. H. Oh, Phys. Rev. Lett. 108, 030402 (2012), DOI 10.1103/PhysRevLett.108.030402.  Construction: published 13-ray seed. The audit verifies 13 projectively distinct rays and 4 complete contexts.\par\endgroup
\begin{longtable}{@{}r l@{}}
\toprule
ray & representative\\
\midrule
\endfirsthead
\toprule
ray & representative\\
\midrule
\endhead
$v_{1}$ & $(1,0,0)$\\
$v_{2}$ & $(0,1,0)$\\
$v_{3}$ & $(0,0,1)$\\
$v_{4}$ & $(0,1,1)$\\
$v_{5}$ & $(0,1,-1)$\\
$v_{6}$ & $(1,0,1)$\\
$v_{7}$ & $(1,0,-1)$\\
$v_{8}$ & $(1,1,0)$\\
$v_{9}$ & $(1,-1,0)$\\
$v_{10}$ & $(1,1,1)$\\
$v_{11}$ & $(1,1,-1)$\\
$v_{12}$ & $(1,-1,1)$\\
$v_{13}$ & $(1,-1,-1)$\\
\bottomrule
\end{longtable}
\begin{longtable}{@{}r l@{}}
\toprule
context & mutually orthogonal rays\\
\midrule
\endfirsthead
\toprule
context & mutually orthogonal rays\\
\midrule
\endhead
$C_{1}$ & $\{v_{1},v_{2},v_{3}\}$\\
$C_{2}$ & $\{v_{1},v_{4},v_{5}\}$\\
$C_{3}$ & $\{v_{2},v_{6},v_{7}\}$\\
$C_{4}$ & $\{v_{3},v_{8},v_{9}\}$\\
\bottomrule
\end{longtable}
\subsection{Yu--Oh 25-ray orthogonal completion}
\begingroup\sloppy Source: S. Yu and C. H. Oh, Phys. Rev. Lett. 108, 030402 (2012), DOI 10.1103/PhysRevLett.108.030402.  Construction: cross-product completion of every orthogonal pair in the 13-ray seed. The audit verifies 25 projectively distinct rays and 16 complete contexts.\par\endgroup
\begin{longtable}{@{}r l@{}}
\toprule
ray & representative\\
\midrule
\endfirsthead
\toprule
ray & representative\\
\midrule
\endhead
$v_{1}$ & $(0,0,1)$\\
$v_{2}$ & $(0,1,-1)$\\
$v_{3}$ & $(0,1,0)$\\
$v_{4}$ & $(0,1,1)$\\
$v_{5}$ & $(1,-1,-1)$\\
$v_{6}$ & $(1,-1,-2)$\\
$v_{7}$ & $(1,-1,0)$\\
$v_{8}$ & $(1,-1,1)$\\
$v_{9}$ & $(1,-1,2)$\\
$v_{10}$ & $(1,-2,-1)$\\
$v_{11}$ & $(1,-2,1)$\\
$v_{12}$ & $(1,0,-1)$\\
$v_{13}$ & $(1,0,0)$\\
$v_{14}$ & $(1,0,1)$\\
$v_{15}$ & $(1,1,-1)$\\
$v_{16}$ & $(1,1,-2)$\\
$v_{17}$ & $(1,1,0)$\\
$v_{18}$ & $(1,1,1)$\\
$v_{19}$ & $(1,1,2)$\\
$v_{20}$ & $(1,2,-1)$\\
$v_{21}$ & $(1,2,1)$\\
$v_{22}$ & $(2,-1,-1)$\\
$v_{23}$ & $(2,-1,1)$\\
$v_{24}$ & $(2,1,-1)$\\
$v_{25}$ & $(2,1,1)$\\
\bottomrule
\end{longtable}
\begin{longtable}{@{}r l@{}}
\toprule
context & mutually orthogonal rays\\
\midrule
\endfirsthead
\toprule
context & mutually orthogonal rays\\
\midrule
\endhead
$C_{1}$ & $\{v_{1},v_{3},v_{13}\}$\\
$C_{2}$ & $\{v_{1},v_{7},v_{17}\}$\\
$C_{3}$ & $\{v_{2},v_{4},v_{13}\}$\\
$C_{4}$ & $\{v_{2},v_{5},v_{25}\}$\\
$C_{5}$ & $\{v_{2},v_{18},v_{22}\}$\\
$C_{6}$ & $\{v_{3},v_{12},v_{14}\}$\\
$C_{7}$ & $\{v_{4},v_{8},v_{24}\}$\\
$C_{8}$ & $\{v_{4},v_{15},v_{23}\}$\\
$C_{9}$ & $\{v_{5},v_{9},v_{17}\}$\\
$C_{10}$ & $\{v_{5},v_{14},v_{20}\}$\\
$C_{11}$ & $\{v_{6},v_{8},v_{17}\}$\\
$C_{12}$ & $\{v_{7},v_{15},v_{19}\}$\\
$C_{13}$ & $\{v_{7},v_{16},v_{18}\}$\\
$C_{14}$ & $\{v_{8},v_{12},v_{21}\}$\\
$C_{15}$ & $\{v_{10},v_{14},v_{15}\}$\\
$C_{16}$ & $\{v_{11},v_{12},v_{18}\}$\\
\bottomrule
\end{longtable}
\subsection{Specker bug}
\begingroup\sloppy Source: J. Tkadlec, Greechie Diagrams of Small Quantum Logics with Small State Spaces, Int. J. Theor. Phys. 37, 203--209 (1998), Fig. 1a, DOI 10.1023/A:1026646229896; A. Cabello Quintero, Pruebas algebraicas de imposibilidad de variables ocultas en mecanica cuantica, PhD thesis, Universidad Complutense de Madrid (1996), Fig. 2.5 and printed p. 44.  Construction: coordinate-labelled 13-line realization; Cabello supplies the parametric B-11/B-13 family. The audit verifies 13 projectively distinct rays and 7 complete contexts.\par\endgroup
\begin{longtable}{@{}r l@{}}
\toprule
ray & representative\\
\midrule
\endfirsthead
\toprule
ray & representative\\
\midrule
\endhead
$v_{1}$ & $(1,\sqrt{2},1)$\\
$v_{2}$ & $(1,0,-1)$\\
$v_{3}$ & $(1,-\sqrt{2},1)$\\
$v_{4}$ & $(1,\sqrt{2},-1)$\\
$v_{5}$ & $(1,0,1)$\\
$v_{6}$ & $(1,-\sqrt{2},-1)$\\
$v_{7}$ & $(1,-1/2\,\sqrt{2},0)$\\
$v_{8}$ & $(1,1/2\,\sqrt{2},0)$\\
$v_{9}$ & $(0,1,0)$\\
$v_{10}$ & $(1,\sqrt{2},-3)$\\
$v_{11}$ & $(1,\sqrt{2},3)$\\
$v_{12}$ & $(1,-\sqrt{2},-3)$\\
$v_{13}$ & $(1,-\sqrt{2},3)$\\
\bottomrule
\end{longtable}
\begin{longtable}{@{}r l@{}}
\toprule
context & mutually orthogonal rays\\
\midrule
\endfirsthead
\toprule
context & mutually orthogonal rays\\
\midrule
\endhead
$C_{1}$ & $\{v_{1},v_{2},v_{3}\}$\\
$C_{2}$ & $\{v_{1},v_{7},v_{10}\}$\\
$C_{3}$ & $\{v_{2},v_{5},v_{9}\}$\\
$C_{4}$ & $\{v_{3},v_{8},v_{12}\}$\\
$C_{5}$ & $\{v_{4},v_{5},v_{6}\}$\\
$C_{6}$ & $\{v_{4},v_{7},v_{11}\}$\\
$C_{7}$ & $\{v_{6},v_{8},v_{13}\}$\\
\bottomrule
\end{longtable}
\subsection{Kochen--Specker \texorpdfstring{$\Gamma_3$}{Gamma3}}
\begingroup\sloppy Source: J. Tkadlec, Greechie Diagrams of Small Quantum Logics with Small State Spaces, Int. J. Theor. Phys. 37, 203--209 (1998), Fig. 1b, DOI 10.1023/A:1026646229896.  Construction: orthogonal completion of the coordinate-labelled displayed rays. The audit verifies 27 projectively distinct rays and 17 complete contexts.\par\endgroup
\begin{longtable}{@{}r l@{}}
\toprule
ray & representative\\
\midrule
\endfirsthead
\toprule
ray & representative\\
\midrule
\endhead
$v_{1}$ & $(0,0,1)$\\
$v_{2}$ & $(0,1,-1)$\\
$v_{3}$ & $(0,1,0)$\\
$v_{4}$ & $(0,1,1)$\\
$v_{5}$ & $(1,-1/2\,\sqrt{2},-1/2\,\sqrt{2})$\\
$v_{6}$ & $(1,-1/2\,\sqrt{2},-3/2\,\sqrt{2})$\\
$v_{7}$ & $(1,-1/2\,\sqrt{2},0)$\\
$v_{8}$ & $(1,-1/2\,\sqrt{2},1/2\,\sqrt{2})$\\
$v_{9}$ & $(1,-1/2\,\sqrt{2},3/2\,\sqrt{2})$\\
$v_{10}$ & $(1,-\sqrt{2},-1)$\\
$v_{11}$ & $(1,-\sqrt{2},-3)$\\
$v_{12}$ & $(1,-\sqrt{2},0)$\\
$v_{13}$ & $(1,-\sqrt{2},1)$\\
$v_{14}$ & $(1,-\sqrt{2},3)$\\
$v_{15}$ & $(1,0,-1)$\\
$v_{16}$ & $(1,0,0)$\\
$v_{17}$ & $(1,0,1)$\\
$v_{18}$ & $(1,1/2\,\sqrt{2},-1/2\,\sqrt{2})$\\
$v_{19}$ & $(1,1/2\,\sqrt{2},-3/2\,\sqrt{2})$\\
$v_{20}$ & $(1,1/2\,\sqrt{2},0)$\\
$v_{21}$ & $(1,1/2\,\sqrt{2},1/2\,\sqrt{2})$\\
$v_{22}$ & $(1,1/2\,\sqrt{2},3/2\,\sqrt{2})$\\
$v_{23}$ & $(1,\sqrt{2},-1)$\\
$v_{24}$ & $(1,\sqrt{2},-3)$\\
$v_{25}$ & $(1,\sqrt{2},0)$\\
$v_{26}$ & $(1,\sqrt{2},1)$\\
$v_{27}$ & $(1,\sqrt{2},3)$\\
\bottomrule
\end{longtable}
\begin{longtable}{@{}r l@{}}
\toprule
context & mutually orthogonal rays\\
\midrule
\endfirsthead
\toprule
context & mutually orthogonal rays\\
\midrule
\endhead
$C_{1}$ & $\{v_{1},v_{3},v_{16}\}$\\
$C_{2}$ & $\{v_{1},v_{7},v_{25}\}$\\
$C_{3}$ & $\{v_{1},v_{12},v_{20}\}$\\
$C_{4}$ & $\{v_{2},v_{4},v_{16}\}$\\
$C_{5}$ & $\{v_{2},v_{5},v_{21}\}$\\
$C_{6}$ & $\{v_{3},v_{15},v_{17}\}$\\
$C_{7}$ & $\{v_{4},v_{8},v_{18}\}$\\
$C_{8}$ & $\{v_{5},v_{9},v_{25}\}$\\
$C_{9}$ & $\{v_{6},v_{8},v_{25}\}$\\
$C_{10}$ & $\{v_{7},v_{23},v_{27}\}$\\
$C_{11}$ & $\{v_{7},v_{24},v_{26}\}$\\
$C_{12}$ & $\{v_{10},v_{14},v_{20}\}$\\
$C_{13}$ & $\{v_{10},v_{17},v_{23}\}$\\
$C_{14}$ & $\{v_{11},v_{13},v_{20}\}$\\
$C_{15}$ & $\{v_{12},v_{18},v_{22}\}$\\
$C_{16}$ & $\{v_{12},v_{19},v_{21}\}$\\
$C_{17}$ & $\{v_{13},v_{15},v_{26}\}$\\
\bottomrule
\end{longtable}
\subsection{Tkadlec nonunital configuration}
\begingroup\sloppy Source: J. Tkadlec, Greechie Diagrams of Small Quantum Logics with Small State Spaces, Int. J. Theor. Phys. 37, 203--209 (1998), Fig. 2, DOI 10.1023/A:1026646229896.  Construction: orthogonal completion of the 25 rays marked in Fig. 2. The audit verifies 37 projectively distinct rays and 26 complete contexts.\par\endgroup
\begin{longtable}{@{}r l@{}}
\toprule
ray & representative\\
\midrule
\endfirsthead
\toprule
ray & representative\\
\midrule
\endhead
$v_{1}$ & $(0,0,1)$\\
$v_{2}$ & $(0,1,-1)$\\
$v_{3}$ & $(0,1,0)$\\
$v_{4}$ & $(0,1,1)$\\
$v_{5}$ & $(1,-1,-1)$\\
$v_{6}$ & $(1,-1,-2)$\\
$v_{7}$ & $(1,-1,0)$\\
$v_{8}$ & $(1,-1,1)$\\
$v_{9}$ & $(1,-1,2)$\\
$v_{10}$ & $(1,-2,-1)$\\
$v_{11}$ & $(1,-2,1)$\\
$v_{12}$ & $(1,-5,-2)$\\
$v_{13}$ & $(1,-5,2)$\\
$v_{14}$ & $(1,0,-1)$\\
$v_{15}$ & $(1,0,-2)$\\
$v_{16}$ & $(1,0,0)$\\
$v_{17}$ & $(1,0,1)$\\
$v_{18}$ & $(1,0,2)$\\
$v_{19}$ & $(1,1,-1)$\\
$v_{20}$ & $(1,1,-2)$\\
$v_{21}$ & $(1,1,0)$\\
$v_{22}$ & $(1,1,1)$\\
$v_{23}$ & $(1,1,2)$\\
$v_{24}$ & $(1,2,-1)$\\
$v_{25}$ & $(1,2,1)$\\
$v_{26}$ & $(1,5,-2)$\\
$v_{27}$ & $(1,5,2)$\\
$v_{28}$ & $(2,-1,-1)$\\
$v_{29}$ & $(2,-1,1)$\\
$v_{30}$ & $(2,-5,-1)$\\
$v_{31}$ & $(2,-5,1)$\\
$v_{32}$ & $(2,0,-1)$\\
$v_{33}$ & $(2,0,1)$\\
$v_{34}$ & $(2,1,-1)$\\
$v_{35}$ & $(2,1,1)$\\
$v_{36}$ & $(2,5,-1)$\\
$v_{37}$ & $(2,5,1)$\\
\bottomrule
\end{longtable}
\begin{longtable}{@{}r l@{}}
\toprule
context & mutually orthogonal rays\\
\midrule
\endfirsthead
\toprule
context & mutually orthogonal rays\\
\midrule
\endhead
$C_{1}$ & $\{v_{1},v_{3},v_{16}\}$\\
$C_{2}$ & $\{v_{1},v_{7},v_{21}\}$\\
$C_{3}$ & $\{v_{2},v_{4},v_{16}\}$\\
$C_{4}$ & $\{v_{2},v_{5},v_{35}\}$\\
$C_{5}$ & $\{v_{2},v_{22},v_{28}\}$\\
$C_{6}$ & $\{v_{3},v_{14},v_{17}\}$\\
$C_{7}$ & $\{v_{3},v_{15},v_{33}\}$\\
$C_{8}$ & $\{v_{3},v_{18},v_{32}\}$\\
$C_{9}$ & $\{v_{4},v_{8},v_{34}\}$\\
$C_{10}$ & $\{v_{4},v_{19},v_{29}\}$\\
$C_{11}$ & $\{v_{5},v_{9},v_{21}\}$\\
$C_{12}$ & $\{v_{5},v_{17},v_{24}\}$\\
$C_{13}$ & $\{v_{6},v_{8},v_{21}\}$\\
$C_{14}$ & $\{v_{6},v_{26},v_{33}\}$\\
$C_{15}$ & $\{v_{7},v_{19},v_{23}\}$\\
$C_{16}$ & $\{v_{7},v_{20},v_{22}\}$\\
$C_{17}$ & $\{v_{8},v_{14},v_{25}\}$\\
$C_{18}$ & $\{v_{9},v_{27},v_{32}\}$\\
$C_{19}$ & $\{v_{10},v_{17},v_{19}\}$\\
$C_{20}$ & $\{v_{11},v_{14},v_{22}\}$\\
$C_{21}$ & $\{v_{12},v_{20},v_{33}\}$\\
$C_{22}$ & $\{v_{13},v_{23},v_{32}\}$\\
$C_{23}$ & $\{v_{15},v_{29},v_{37}\}$\\
$C_{24}$ & $\{v_{15},v_{31},v_{35}\}$\\
$C_{25}$ & $\{v_{18},v_{28},v_{36}\}$\\
$C_{26}$ & $\{v_{18},v_{30},v_{34}\}$\\
\bottomrule
\end{longtable}
\subsection{Cabello 18--9}
\begingroup\sloppy Source: A. Cabello, J. M. Estebaranz, and G. Garcia-Alcaine, Phys. Lett. A 212, 183--187 (1996), DOI 10.1016/0375-9601(96)00134-X.  Construction: published real 18-ray coordinatization. The audit verifies 18 projectively distinct rays and 9 complete contexts.\par\endgroup
\begin{longtable}{@{}r l@{}}
\toprule
ray & representative\\
\midrule
\endfirsthead
\toprule
ray & representative\\
\midrule
\endhead
$v_{1}$ & $(1,0,0,0)$\\
$v_{2}$ & $(0,0,1,-1)$\\
$v_{3}$ & $(0,0,1,1)$\\
$v_{4}$ & $(0,1,0,0)$\\
$v_{5}$ & $(0,1,-1,0)$\\
$v_{6}$ & $(0,0,0,1)$\\
$v_{7}$ & $(0,1,1,0)$\\
$v_{8}$ & $(1,-1,-1,-1)$\\
$v_{9}$ & $(1,1,1,-1)$\\
$v_{10}$ & $(1,0,0,1)$\\
$v_{11}$ & $(0,1,0,-1)$\\
$v_{12}$ & $(1,1,-1,1)$\\
$v_{13}$ & $(1,0,1,0)$\\
$v_{14}$ & $(1,1,1,1)$\\
$v_{15}$ & $(1,0,-1,0)$\\
$v_{16}$ & $(1,-1,1,-1)$\\
$v_{17}$ & $(1,-1,0,0)$\\
$v_{18}$ & $(1,1,-1,-1)$\\
\bottomrule
\end{longtable}
\begin{longtable}{@{}r l@{}}
\toprule
context & mutually orthogonal rays\\
\midrule
\endfirsthead
\toprule
context & mutually orthogonal rays\\
\midrule
\endhead
$C_{1}$ & $\{v_{1},v_{2},v_{3},v_{4}\}$\\
$C_{2}$ & $\{v_{1},v_{5},v_{6},v_{7}\}$\\
$C_{3}$ & $\{v_{2},v_{14},v_{17},v_{18}\}$\\
$C_{4}$ & $\{v_{3},v_{9},v_{12},v_{17}\}$\\
$C_{5}$ & $\{v_{4},v_{6},v_{13},v_{15}\}$\\
$C_{6}$ & $\{v_{5},v_{8},v_{9},v_{10}\}$\\
$C_{7}$ & $\{v_{7},v_{10},v_{16},v_{18}\}$\\
$C_{8}$ & $\{v_{8},v_{11},v_{12},v_{13}\}$\\
$C_{9}$ & $\{v_{11},v_{14},v_{15},v_{16}\}$\\
\bottomrule
\end{longtable}
\subsection{Peres--Mermin/Peres 24}
\begingroup\sloppy Source: D. N. Mermin, Phys. Rev. Lett. 65, 3373--3377 (1990), DOI 10.1103/PhysRevLett.65.3373; A. Peres, J. Phys. A 24, L175--L178 (1991), DOI 10.1088/0305-4470/24/4/003.  Construction: exact simultaneous diagonalization of the six Peres--Mermin contexts. The audit verifies 24 projectively distinct rays and 24 complete contexts.\par\endgroup
\begin{longtable}{@{}r l@{}}
\toprule
ray & representative\\
\midrule
\endfirsthead
\toprule
ray & representative\\
\midrule
\endhead
$v_{1}$ & $(0,0,0,1)$\\
$v_{2}$ & $(0,0,1,0)$\\
$v_{3}$ & $(0,1,0,0)$\\
$v_{4}$ & $(1,0,0,0)$\\
$v_{5}$ & $(1,-1,-1,1)$\\
$v_{6}$ & $(1,-1,1,-1)$\\
$v_{7}$ & $(1,1,-1,-1)$\\
$v_{8}$ & $(1,1,1,1)$\\
$v_{9}$ & $(1,-1,-1,-1)$\\
$v_{10}$ & $(1,-1,1,1)$\\
$v_{11}$ & $(1,1,-1,1)$\\
$v_{12}$ & $(1,1,1,-1)$\\
$v_{13}$ & $(0,0,1,-1)$\\
$v_{14}$ & $(0,0,1,1)$\\
$v_{15}$ & $(1,-1,0,0)$\\
$v_{16}$ & $(1,1,0,0)$\\
$v_{17}$ & $(0,1,0,-1)$\\
$v_{18}$ & $(0,1,0,1)$\\
$v_{19}$ & $(1,0,-1,0)$\\
$v_{20}$ & $(1,0,1,0)$\\
$v_{21}$ & $(0,1,-1,0)$\\
$v_{22}$ & $(0,1,1,0)$\\
$v_{23}$ & $(1,0,0,-1)$\\
$v_{24}$ & $(1,0,0,1)$\\
\bottomrule
\end{longtable}
\begin{longtable}{@{}r l@{}}
\toprule
context & mutually orthogonal rays\\
\midrule
\endfirsthead
\toprule
context & mutually orthogonal rays\\
\midrule
\endhead
$C_{1}$ & $\{v_{1},v_{2},v_{3},v_{4}\}$\\
$C_{2}$ & $\{v_{1},v_{2},v_{15},v_{16}\}$\\
$C_{3}$ & $\{v_{1},v_{3},v_{19},v_{20}\}$\\
$C_{4}$ & $\{v_{1},v_{4},v_{21},v_{22}\}$\\
$C_{5}$ & $\{v_{2},v_{3},v_{23},v_{24}\}$\\
$C_{6}$ & $\{v_{2},v_{4},v_{17},v_{18}\}$\\
$C_{7}$ & $\{v_{3},v_{4},v_{13},v_{14}\}$\\
$C_{8}$ & $\{v_{5},v_{6},v_{7},v_{8}\}$\\
$C_{9}$ & $\{v_{5},v_{6},v_{14},v_{16}\}$\\
$C_{10}$ & $\{v_{5},v_{7},v_{18},v_{20}\}$\\
$C_{11}$ & $\{v_{5},v_{8},v_{21},v_{23}\}$\\
$C_{12}$ & $\{v_{6},v_{7},v_{22},v_{24}\}$\\
$C_{13}$ & $\{v_{6},v_{8},v_{17},v_{19}\}$\\
$C_{14}$ & $\{v_{7},v_{8},v_{13},v_{15}\}$\\
$C_{15}$ & $\{v_{9},v_{10},v_{11},v_{12}\}$\\
$C_{16}$ & $\{v_{9},v_{10},v_{13},v_{16}\}$\\
$C_{17}$ & $\{v_{9},v_{11},v_{17},v_{20}\}$\\
$C_{18}$ & $\{v_{9},v_{12},v_{21},v_{24}\}$\\
$C_{19}$ & $\{v_{10},v_{11},v_{22},v_{23}\}$\\
$C_{20}$ & $\{v_{10},v_{12},v_{18},v_{19}\}$\\
$C_{21}$ & $\{v_{11},v_{12},v_{14},v_{15}\}$\\
$C_{22}$ & $\{v_{13},v_{14},v_{15},v_{16}\}$\\
$C_{23}$ & $\{v_{17},v_{18},v_{19},v_{20}\}$\\
$C_{24}$ & $\{v_{21},v_{22},v_{23},v_{24}\}$\\
\bottomrule
\end{longtable}
\subsection{Completed KCBS pentagon}
Put \(h=\sqrt{\cos(\pi/5)}\).  The ten ray representatives used by the numerical branch are:
\begin{longtable}{@{}r l@{}}
\toprule ray & representative\\ \midrule
\endfirsthead \toprule ray & representative\\ \midrule \endhead
$p_{0}$ & $(1,0,h)$\\
$p_{1}$ & $(\cos(4\pi/5),\sin(4\pi/5),h)$\\
$p_{2}$ & $(\cos(8\pi/5),\sin(8\pi/5),h)$\\
$p_{3}$ & $(\cos(12\pi/5),\sin(12\pi/5),h)$\\
$p_{4}$ & $(\cos(16\pi/5),\sin(16\pi/5),h)$\\
$q_{0}=p_{0}\mathbin{\times}p_{1}$ & $(h[\sin(0)-\sin(4\pi/5)],h[\cos(4\pi/5)-\cos(0)],\sin(4\pi/5-0))$\\
$q_{1}=p_{1}\mathbin{\times}p_{2}$ & $(h[\sin(4\pi/5)-\sin(8\pi/5)],h[\cos(8\pi/5)-\cos(4\pi/5)],\sin(8\pi/5-4\pi/5))$\\
$q_{2}=p_{2}\mathbin{\times}p_{3}$ & $(h[\sin(8\pi/5)-\sin(12\pi/5)],h[\cos(12\pi/5)-\cos(8\pi/5)],\sin(12\pi/5-8\pi/5))$\\
$q_{3}=p_{3}\mathbin{\times}p_{4}$ & $(h[\sin(12\pi/5)-\sin(16\pi/5)],h[\cos(16\pi/5)-\cos(12\pi/5)],\sin(16\pi/5-12\pi/5))$\\
$q_{4}=p_{4}\mathbin{\times}p_{0}$ & $(h[\sin(16\pi/5)-\sin(0)],h[\cos(0)-\cos(16\pi/5)],\sin(0-16\pi/5))$\\
\bottomrule
\end{longtable}
The five contexts are
\[
\{p_0,q_0,p_1\},\quad \{p_1,q_1,p_2\},\quad \{p_2,q_2,p_3\},
\quad \{p_3,q_3,p_4\},\quad \{p_4,q_4,p_0\}.
\]
